\section{Consolidation Question Q1: Novelty and Equivalence Audit}
\label{sec:consolidation-q1}

The first consolidation question asks whether the frozen stateful event encoder is mathematically distinct from the closest existing communication-efficient optimization and federated-learning methods. This question must be answered before another mechanism is added. A method can be empirically useful without being algorithmically new, and Experiment~9 already established exact equivalences for two earlier recurrences. The present audit therefore compares the \emph{complete} frozen V1 mechanism against the nearest prior constructions in common notation.

\subsection{Frozen V1 in compressor notation}

Let $u_i^r\in\mathbb{R}^d$ denote the update signal generated by client $i$. Frozen V1 maintains a client-side state
\begin{equation}
 z_i^{r+1}=\rho z_i^r+\alpha_i^r u_i^r,
 \label{eq:q1-v1-state}
\end{equation}
and generates a coordinate event whenever
\begin{equation}
 |z_{ij}^{r+1}|\geq \Delta_j.
\end{equation}
A fired coordinate transmits only
\begin{equation}
 (j,s_{ij}),\qquad s_{ij}=\operatorname{sign}(z_{ij}^{r+1}),
\end{equation}
then undergoes full reset,
\begin{equation}
 z_{ij}^{r+1}\leftarrow0,
 \label{eq:q1-v1-reset}
\end{equation}
and the server applies
\begin{equation}
 w_j^+=w_j+q_r s_{ij}.
 \label{eq:q1-v1-jump}
\end{equation}
The defining point for the novelty audit is that the communication threshold $\Delta_j$ and the model-update quantum $q_r$ are independent design variables. In the successful neural-network experiments, $q_r$ is substantially smaller than $\Delta_j$ and is annealed more strongly than the evidence memory. Consequently, the transmitted model update is not a reconstruction of the accumulator value and the accumulator overshoot is intentionally discarded by the full reset.

\subsection{Previously established exact equivalences}

Experiment~9 already excludes two broad novelty claims. For subtractive reset,
\begin{equation}
 z^{r+1}=\rho z^r+u^r-\Delta s^r,
\end{equation}
rescaling the state gives a decayed error-feedback recurrence. Therefore subtractive FedLIF cannot be claimed as a new residual-feedback mechanism. Error feedback and memory-decayed error feedback are established techniques in distributed optimization and communication-efficient learning~\cite{stich2018memory,chen2022co3}.

For full reset, the local inter-event recurrence is exactly a resetting exponential moving average after rescaling. If
\begin{equation}
 z^{r+1}=\rho z^r-c g^r,
\end{equation}
and
\begin{equation}
 m^r=-\frac{1-\rho}{c}z^r,
\end{equation}
then between events
\begin{equation}
 m^{r+1}=\rho m^r+(1-\rho)g^r,
\end{equation}
with an equivalent threshold and reset $m^+=0$. Hence the local filtering recurrence itself is also not novel. The only remaining question is whether the \emph{complete communication operator} in \eqref{eq:q1-v1-state}--\eqref{eq:q1-v1-jump} already exists.

\subsection{The closest predecessor: Strom's gradient-residual compressor}

The strongest result of the present literature audit is that the nearest predecessor is considerably closer than generic signSGD or Top-$K$ sparsification. Strom's 2015 distributed-SGD method~\cite{strom2015distributed} maintains a gradient residual on every worker. After each mini-batch the stochastic gradient is accumulated into that residual. When coordinate $j$ exceeds a fixed threshold $\tau$, the worker communicates a signed quantum $+\tau$ or $-\tau$ together with the coordinate address and subtracts the same quantum from the residual. In simplified notation,
\begin{align}
 r_i^{r+1} &= r_i^r+u_i^r,\\
 r_{ij}^{r+1}>\tau
 &\Rightarrow
 \begin{cases}
   \text{send }(j,+\tau),\\
   r_{ij}^{r+1}\leftarrow r_{ij}^{r+1}-\tau,
 \end{cases}\\
 r_{ij}^{r+1}<-\tau
 &\Rightarrow
 \begin{cases}
   \text{send }(j,-\tau),\\
   r_{ij}^{r+1}\leftarrow r_{ij}^{r+1}+\tau.
 \end{cases}
\end{align}
The method also explicitly supports asynchronous operation. Strom further observes that one-bit quantization is sufficient and therefore communicates signed threshold-sized weight deltas indexed by coordinate.

This prior work establishes that the broad mechanism cannot be claimed as new.
\begin{takeawaybox}[title={Q1 prior-art boundary}]
Per-coordinate temporal gradient accumulation, threshold crossing, address/sign communication, and asynchronous sparse model updates already exist in closely related residual-pulse methods. The novelty question must therefore concern the complete V1 operator, not these ingredients individually.
\end{takeawaybox}
It also gives a concrete algorithmic interpretation to the equivalence found in Experiment~9. In the special case $\rho=1$, $q=\Delta=\tau$, and subtractive reset, the earlier subtractive FedLIF mechanism is essentially Strom-style residual compression.

Frozen V1 is nevertheless not algebraically identical to Strom's method. There are three structural differences:
\begin{enumerate}
    \item V1 uses finite-memory/leaky accumulation, $\rho<1$, rather than an indefinitely conserved gradient residual.
    \item V1 performs \emph{full reset} after a threshold crossing rather than subtracting the transmitted quantum and retaining the remaining residual.
    \item Most importantly, V1 decouples the model quantum $q_r$ from the trigger threshold $\Delta$. Strom's one-bit update has magnitude $\tau$, the same quantity that defines the residual threshold. In V1, $q_r$ can be much smaller than $\Delta$ and can follow an independent annealing law.
\end{enumerate}
The empirical value of these differences must be demonstrated directly. Their mere presence is not sufficient to support a strong novelty claim.

\subsection{Other close prior families}

\paragraph{LENA: self-triggered accumulated-error uploads.}
LENA~\cite{ghadikolaei2021lena} maintains a worker error vector, accumulates the discrepancy between locally computed gradients and server-applied updates, and self-triggers communication through a significance filter. When upload occurs, the memory can be cleared completely. This is close to the ``accumulate--trigger--clear'' logic of V1. The key differences are that LENA uses a vector-level significance test, uploads the accumulated vector rather than a sign/address pulse, and lets the server use a drift estimate when the worker remains silent. Thus LENA establishes that self-triggered accumulated memory with full clearing is prior art, but it does not reproduce V1's coordinate pulse encoder.

\paragraph{Gradient residual sparsification and ternary compression.}
Stich et al., Deep Gradient Compression, AdaComp, Sparse Binary Compression, Sparse Ternary Compression, and ASTC all combine some form of residual/memory accumulation with sparse transmission~\cite{stich2018memory,lin2017dgc,chen2018adacomp,sattler2019sbc,sattler2019stc,mao2022astc}. ASTC is particularly relevant because its nodes accumulate gradient information, threshold sparse coordinates, and ternarize selected values. These methods generally preserve or compensate the compression error. Their sparse-selection rule is also typically a Top-$K$, adaptive sparsity, entropy, or residual criterion rather than the V1 first-passage/full-reset rule. Sparse Ternary Compression is especially important for the later FedAvg audit because it was designed explicitly for federated learning and compresses both directions.

\paragraph{One-bit and sign methods.}
1-bit SGD and signSGD establish that one-bit directional information can support distributed optimization~\cite{seide2014onebit,bernstein2018signsgd}. They differ fundamentally in temporal structure. signSGD communicates signs at every optimization iteration and has no persistent coordinate state or silent first-passage process. 1-bit SGD retains quantization error, placing it closer to error feedback than to frozen full-reset V1.

\paragraph{Event-triggered and lazy distributed learning.}
LAG and subsequent event-triggered learning methods reduce communication by transmitting only when a gradient, model, or update change is sufficiently important~\cite{chen2018lag,yoon2025etfl}. Their trigger is generally applied to a vector-level model/update discrepancy and their messages retain model/update magnitude. TLAQC combines adaptive communication sparsification with accumulated quantization errors and retained model deltas~\cite{ren2023tlaqc}. These works establish that event-triggered FL itself is not a novelty claim available to V1.

\paragraph{Approximate synchronous parallelism.}
Gaia's Approximate Synchronous Parallel synchronization model suppresses insignificant parameter updates and allows small updates to accumulate before communication~\cite{hsieh2017gaia}. This reinforces the conclusion that ``accumulate until significant'' is a long-established systems idea. Gaia communicates the accumulated model update itself and does not implement V1's sign-only fixed-quantum reset encoder.

\paragraph{Level-crossing distributed estimation.}
Utlu, Kilic, and Kozat~\cite{utlu2017level} provide the closest non-learning communication analogy. Their distributed estimator communicates parameter information only when a known quantization level is crossed. For a single level crossing, only the direction needs to be transmitted and the receiver reconstructs the crossed level. This establishes that event timing plus a direction bit can encode the evolution of a distributed parameter trajectory. The construction is nevertheless different from V1: the level crossing is performed on a local parameter estimate relative to the receiver's last reconstructed level, not on a leaky accumulated learning-evidence state, and the reconstruction increment is tied to the quantizer's level spacing.

\subsection{Why frozen V1 is not exact error feedback}

The distinction from residual-conserving compressors can be written as an invariant. Let $v_j$ denote the pre-compression state of a fired coordinate. A conventional residual/error-feedback compressor preserves
\begin{equation}
 e_j^+=v_j-C(v_j),
 \label{eq:q1-ef-residual}
\end{equation}
where $C(v_j)$ is the communicated update. Strom's one-bit special case has $C(v_j)=\Delta\operatorname{sign}(v_j)$ and therefore retains the overshoot $v_j-C(v_j)$.

Frozen V1 instead applies
\begin{equation}
 C_r(v_j)=q_r\operatorname{sign}(v_j),
 \qquad
 z_j^+=0.
 \label{eq:q1-v1-lossy}
\end{equation}
For \eqref{eq:q1-v1-lossy} to be an exact instance of \eqref{eq:q1-ef-residual}, every fired coordinate would need to satisfy
\begin{equation}
 v_j=q_r\operatorname{sign}(v_j).
\end{equation}
That condition does not hold in V1. The trigger only guarantees $|v_j|\geq\Delta$, while successful experiments deliberately use $q_r<\Delta$ and subsequently decrease $q_r$. Hence the fired-state overshoot and much of its magnitude are intentionally discarded. Magnitude is represented statistically through repeated event occurrence rather than conserved exactly in a residual.

This lossy reset is not automatically beneficial. It is precisely the point that the nearest-neighbor empirical audit must test.

\subsection{Structural comparison}

\begin{table}[ht]
\centering
\small
\caption{Structural comparison of frozen V1 with its closest identified predecessors. ``Coord.'' denotes a coordinate-level operation.}
\label{tab:q1-structure}
\begin{adjustbox}{max width=\textwidth}
\begin{tabular}{p{3.6cm}ccccc}
\toprule
Property & V1 & Strom & LENA & STC & Level crossing\\
\midrule
Persistent local memory & yes & yes & yes & yes & yes\\
Memory is leaky & yes & no & no & no & no\\
Coordinate threshold crossing & yes & yes & no & no & yes\\
Sign/address-only pulse possible & yes & yes & no & ternary & yes\\
Full memory reset on fired coord. & yes & no & vector clear & no & no\\
Compression error conserved & no & yes & yes & yes & n/a\\
Update quantum tied to threshold & no & yes & n/a & generally no & level spacing\\
Independent annealed model quantum & yes & no & no & no & no\\
Server update on silence & none & none & drift/history & round aggregation & hold level\\
Asynchronous formulation & yes & yes & not primary & round based & yes\\
\bottomrule
\end{tabular}
\end{adjustbox}
\end{table}

Table~\ref{tab:q1-structure} makes the novelty boundary explicit. V1 does not introduce temporal accumulation, sparse threshold communication, sign compression, event-triggered learning, or level-crossing coding individually. Its candidate distinction is the specific lossy pulse encoder used as a federated client-update communication layer.
\begin{takeawaybox}[title={Q1 candidate distinction}]
V1 combines leaky coordinate evidence, first-passage triggering, a sign/address pulse with independently controlled model quantum $q_r$, and full reset of fired evidence. This complete operator---rather than accumulation, thresholding, or sign coding alone---is the object carried into the final FL experiments.
\end{takeawaybox}

\subsection{Analytical Q1 verdict}

The literature audit does not identify an exact prior method with the complete frozen V1 recursion and protocol. However, the nearest prior work is sufficiently close that a broad priority claim would be inappropriate. In particular, Strom already contains the zero-leak, residual-preserving, threshold/sign-pulse sibling of the mechanism, while level-crossing distributed estimation already contains the event-time/direction-code interpretation.

The analytical decision is therefore
\begin{equation}
\boxed{\text{Q1 analytical audit: QUALIFIED PASS}.}
\end{equation}
The final matched campaign supports only the following narrow claim:
\begin{quote}
The proposed communication layer uses a leaky per-coordinate learning-evidence state and converts first-passage events into fixed signed model quanta whose magnitude is decoupled from the communication threshold and annealed independently; fired evidence is fully reset rather than conserved as compression error.
\end{quote}
No manuscript should claim invention of gradient accumulation, error feedback, threshold sparsification, sign/ternary communication, event-triggered FL, or temporal/event coding.

\subsection{P2 literature refresh}

The operator-level search was refreshed on 2026-09-04.  Recent event-triggered
learning methods retain the established model- or gradient-deviation pattern:
DETSGRAD and EventGraD suppress vector communication
\citep{george2020detsgrad,ghosh2022eventgrad}; SPARQ-SGD combines a
model-deviation event with quantization and sparsification
\citep{singh2023sparq}; bidirectional ETFL triggers on differences from the
last communicated models \citep{he2023etfl}; and 2026 event-triggered gossip
uses decentralized local-model deviation \citep{zhai2026gossip}.  These works
do not use the same leaky coordinate evidence, full-clear pulse, and independent
server resolution.

Federated SNN work places LIF dynamics inside the trained network and exchanges
SNN weights, compressed updates, or firing-rate information
\citep{venkatesha2021fedsnn,chaki2023tradeoffs,tao2026sfedhifi}.  Frozen V1
instead wraps conventional FedAvg updates in a LIF-inspired communication
encoder.  Consequently, no claim of ``first neuromorphic FL'' is admissible.
The complete operator comparison and query protocol are maintained in
\texttt{paper/ijcnn2027/RELATED\_WORK.md}.  The search found no exact
predecessor for the full frozen operator, but this is a qualified, dated
absence-search result and must be refreshed before submission.

\subsection{Completed nearest-neighbor experiment}

The bridge experiment requested by Q1 was completed in
Section~\ref{sec:final-baseline-campaign}.  Under matched local learning,
development-only tuning, independent held-out partitions, and symmetric
end-to-end communication accounting, traffic-matched Event-FedAvg outperformed
the residual-conserving Strom sibling on both architectures.  On the MLP,
Event-FedAvg obtained $83.14\%\pm0.14\%$ accuracy at
$183.5\pm2.0$~Mbit, versus $79.51\%\pm0.86\%$ at
$181.0\pm6.7$~Mbit for Strom.  On the compact CNN, it obtained
$81.34\%\pm0.46\%$ at $72.8\pm1.2$~Mbit, versus
$79.88\%\pm1.56\%$ at $82.0\pm1.7$~Mbit.  A quality-selected Strom point
reached slightly better CNN accuracy only at approximately $4.94$ times the
conservative total traffic.

This result strengthens the qualified pass: leakage, full reset, and
decoupled resolution define an empirically meaningful operating point relative
to the closest recurrence.  It still does not establish that each ingredient is
individually necessary, universal dominance, or broad priority over all
event-triggered learning.
